%% slide03_study_note.tex
%% Detailed Study Note --- Slide 3: Research Gaps & Objectives
%% TSA Meeting Preparation | Anoop V Eluvathingal | IIT Madras | 2026
\documentclass[12pt,a4paper]{article}

\usepackage[margin=2.5cm]{geometry}
\usepackage{amsmath,amssymb,bm}
\usepackage{booktabs,tabularx,array,longtable}
\usepackage{graphicx}
\usepackage{xcolor}
\usepackage{tcolorbox}
\usepackage{enumitem}
\usepackage{fancyhdr}
\usepackage{titlesec}
\usepackage{hyperref}
\usepackage{parskip}

%% Custom commands
\newcommand{\kibr}{k_\mathrm{ibr}}
\newcommand{\km}{k_m}
\newcommand{\Zline}{Z_\mathrm{line}}
\newcommand{\Yibr}{Y_\mathrm{ibr}}
\newcommand{\CTI}{\mathrm{CTI}}
\newcommand{\gammath}{\gamma_\mathrm{th}}

%% Colours
\definecolor{iitmnavy}{RGB}{15,40,80}
\definecolor{iitmgold}{RGB}{180,140,0}
\definecolor{lightblue}{RGB}{220,235,250}
\definecolor{lightyellow}{RGB}{255,252,220}
\definecolor{lightred}{RGB}{255,235,235}
\definecolor{lightgreen}{RGB}{230,255,230}
\definecolor{lightpurple}{RGB}{240,230,255}
\definecolor{lightorange}{RGB}{255,243,225}
\definecolor{gapcol}{RGB}{255,230,230}
\definecolor{objcol}{RGB}{220,240,220}

%% tcolorbox environments --- title={#1} prevents pgfkeys parsing commas
\tcbuselibrary{skins,breakable}
\newtcolorbox{keybox}[1]{colback=lightblue,colframe=iitmnavy,
  fonttitle=\bfseries\color{white},title={#1},arc=3pt,boxrule=1pt,breakable}
\newtcolorbox{formulabox}[1]{colback=lightyellow,colframe=iitmgold,
  fonttitle=\bfseries,title={#1},arc=3pt,boxrule=1pt,breakable}
\newtcolorbox{resultbox}[1]{colback=lightgreen,colframe=green!50!black,
  fonttitle=\bfseries,title={#1},arc=3pt,boxrule=1pt,breakable}
\newtcolorbox{warnbox}[1]{colback=lightred,colframe=red!60!black,
  fonttitle=\bfseries,title={#1},arc=3pt,boxrule=1pt,breakable}
\newtcolorbox{archbox}[1]{colback=lightpurple,colframe=iitmnavy!60,
  fonttitle=\bfseries,title={#1},arc=3pt,boxrule=1pt,breakable}
\newtcolorbox{speakbox}{colback=orange!8,colframe=orange!70!black,
  fonttitle=\bfseries,title={What to Say at the TSA Meeting},arc=3pt,
  boxrule=1pt,breakable}

%% Header / footer
\pagestyle{fancy}
\fancyhf{}
\fancyhead[L]{\color{iitmnavy}\small\textbf{TSA Meeting Study Note}}
\fancyhead[R]{\color{iitmnavy}\small Slide 3 --- Research Gaps \& Objectives}
\fancyfoot[C]{\thepage}
\fancyfoot[R]{\small\color{gray}Anoop V Eluvathingal $\cdot$ IIT Madras $\cdot$ 2026}
\renewcommand{\headrulewidth}{1.5pt}
\renewcommand{\headrule}{\color{iitmnavy}\hrule width\headwidth height\headrulewidth}

\titleformat{\section}{\large\bfseries\color{iitmnavy}}{}{0em}{}[\color{iitmnavy}\titlerule]
\titleformat{\subsection}{\normalsize\bfseries\color{iitmnavy!80}}{}{0em}{}

\hypersetup{colorlinks=true,linkcolor=iitmnavy,urlcolor=iitmnavy}

\begin{document}

%% ── Title Block ──────────────────────────────────────────────────────────────
\begin{tcolorbox}[colback=iitmnavy,coltext=white,arc=4pt,boxrule=0pt]
\begin{center}
  {\Large\bfseries TSA Meeting --- Slide 3 Study Note}\\[4pt]
  {\large Research Gaps \& Objectives}\\[6pt]
  {\normalsize 6 Identified Gaps $\cdot$ 6 Research Objectives $\cdot$ Prior Art Positioning}\\[2pt]
  {\small Anoop V Eluvathingal $\cdot$ IIT Madras $\cdot$ PhD Thesis 2026}
\end{center}
\end{tcolorbox}

\vspace{6pt}

%% ── 1. Slide Overview ────────────────────────────────────────────────────────
\section{Slide Overview and Narrative}

\begin{keybox}{What This Slide Shows}
Slide~3 is the \textbf{research positioning} slide.  It answers the two most
important questions a TSA committee will ask: \emph{What has not been done
before?} (identified gaps) and \emph{What exactly are you doing?} (research
objectives).  Every claim of novelty in the thesis is traceable to one of
the six gaps on the left; every experimental chapter in the thesis delivers
one of the six objectives on the right.

\medskip
\textbf{The six gaps} can be grouped into three clusters:
\begin{itemize}[nosep]
  \item \textbf{Standards and modelling gap} (G1, G2):
        No relay standard handles real-time $\kibr$; grid-forming IBR
        protection-control coupling is unstudied.
  \item \textbf{Algorithmic gap} (G3, G4):
        No unified inverse-estimation framework exists; existing digital
        twins lack formal bounded non-blocking update guarantees.
  \item \textbf{System-level gap} (G5, G6):
        WAPC designs ignore IBR-induced coherency breakdown; adversarial
        ML attacks on protection are unmitigated.
\end{itemize}

\textbf{The six objectives} map one-to-one to these clusters:
O1--O2 address the standards/modelling and algorithmic gaps;
O3--O4 address the digital twin and ML gaps; O5--O6 address the
system-level and security gaps.  The slide is deliberately structured
as a symmetrical 6-to-6 map.
\end{keybox}

\subsection{Gap-to-Objective Master Map}

\begin{center}
\small
\begin{tabular}{@{}lp{4.5cm}lp{4.5cm}@{}}
\toprule
\textbf{Gap} & \textbf{Statement} & \textbf{Objective} & \textbf{How It Closes the Gap} \\
\midrule
G1 & No standard accounts for real-time $\kibr$
   & O1: Unified SAMBP LM framework
   & LM identifies $\kibr$ online for all relay types \\[4pt]
G2 & GFM protection-control coupling unstudied
   & O2: Self-adaptive modules (incl.\ microgrid)
   & Microgrid module handles GFL/GFM dual-mode \\[4pt]
G3 & No unified inverse-estimation across relays
   & O1 (unified framework)
   & One LM solver with element-specific model \\[4pt]
G4 & DT lacks bounded non-blocking updates
   & O3: Non-blocking EKF digital twin
   & EKF + bounded update law + non-blocking design \\[4pt]
G5 & WAPC ignores IBR coherency breakdown
   & O5: WAPC with topology-aware IBR coherency
   & WAPC uses $\kibr$ in N-1 coherency checks \\[4pt]
G6 & Adversarial ML attacks unmitigated
   & O4: Physics-informed ML; O6: Cybersecurity layer
   & PINN loss + HMAC GOOSE authentication \\
\bottomrule
\end{tabular}
\end{center}

%% ── 2. The Six Gaps: Deep Dive ───────────────────────────────────────────────
\section{The Six Gaps: Why Each Is Genuinely a Gap}

\subsection{G1 --- No Relay Standard Accounts for Real-Time IBR Admittance $\kibr$}

\begin{warnbox}{Standards Gap: IEC 60255 and IEEE C37 Assume Fixed Fault Current}
The two governing relay standards are:
\begin{itemize}[nosep]
  \item \textbf{IEC~60255-151} (overcurrent relay characteristics): Defines
        ITOC/IDMT curves with fixed pickup $I_p$ and time multiplier TMS.
        No provision for real-time pickup adaptation based on source composition.
  \item \textbf{IEEE~C37.112} (inverse-time relay characteristics): Similarly
        assumes the fault current model is determined at commissioning.
  \item \textbf{IEC~61869} (instrument transformer standards): Specifies CT
        accuracy classes but does not address the distortion of fault current
        by IBR current limiting.
\end{itemize}

The standards were written when the only adjustable parameters were $I_p$,
TMS, and zone reach --- all set offline by a protection engineer.  The concept
of a relay that updates its own settings every 20~ms based on a real-time
$\kibr$ estimate does not appear in any relay standard.

\textbf{Industry gap:} Even the latest ``adaptive relay'' products (Siemens
SIPROTEC~5, ABB REL~670) use ``setting groups'' --- a finite set of pre-configured
offline settings that are switched based on topology events.  They do not compute
relay settings from first-principles model identification during the fault itself.
\end{warnbox}

\subsection{G2 --- GFM Protection-Control Coupling Unstudied}

Grid-forming (GFM) inverters operate as voltage-controlled sources emulating
a virtual synchronous machine.  Unlike GFL inverters, the GFM fault current is
determined by its virtual impedance $Z_v = R_v + \jmath\omega_0 L_v$:
\begin{equation}
  I_\mathrm{sc}^\mathrm{GFM} = \frac{|V_0|}{|Z_v|} \in [1.5, 2.5]\,I_\mathrm{rated}
  \label{eq:gfm_isc}
\end{equation}

The protection challenge: the same $Z_v$ that limits fault current also participates
in the microgrid's voltage regulation and frequency droop.  Changing $Z_v$ to improve
protection performance (e.g., raise fault current for relay visibility) directly
compromises the GFM's control stability.  This \emph{protection-control coupling}
is documented in the power electronics literature (Taul~2020, Paquette~2015) but no
protection relay design has explicitly addressed the constraint.

The SAMBP microgrid module (Ch~5) operates within the GFM control constraints by
adapting \emph{relay settings}, not inverter parameters --- this is the only design
approach that does not require changes to the GFM firmware.

\subsection{G3 --- No Unified Inverse-Estimation Framework Across Relay Types}

Prior work addresses each relay type in isolation:
\begin{itemize}[nosep]
  \item Adamiak~(2018): Adaptive 87L using online slope estimation for HVDC
        lines --- specific to DC sources, not GFL/GFM.
  \item Hooshyar~(2019): Retuned OC settings for high-IBR feeders --- offline
        process, not inverse estimation during fault.
  \item Sadeghi~(2021): Distance relay reach correction for IBR networks ---
        element-specific correction table, no model identification.
\end{itemize}

Each paper proposes a \emph{different} estimation method tailored to \emph{one}
relay type.  There is no published work that:
(a) defines a single mathematical problem class (bounded nonlinear
least-squares) that covers 87L, OC, 87T, and distance simultaneously; and
(b) proves convergence and safety properties for the unified solver.
Chapter~2 fills this gap.

\subsection{G4 --- Digital Twin Approaches Lack Bounded Non-Blocking Update Laws}

Existing power system digital twins (Daniluk~2021, IEEE~PES DT task force~2022)
use one of two approaches:
\begin{itemize}[nosep]
  \item \textbf{Offline batch estimation:} Run optimisation after the event,
        update settings during the next maintenance window.  Latency: hours to days.
        Protection cannot benefit from the estimate in real time.
  \item \textbf{Online blocking estimation:} Run the estimator in the relay's
        main compute loop.  If the estimator takes $>$1~ms, the relay misses its
        trip decision window.  No published protection DT has formal non-blocking
        guarantees.
\end{itemize}

Neither approach provides:
(a) a formal proof that settings stay within safe bounds regardless of estimation
quality, or (b) an asynchronous (non-blocking) design where the estimator runs
in a background thread and the relay always has a valid (if possibly stale)
parameter estimate.  Chapter~9 provides both.

\subsection{G5 --- WAPC Designs Ignore IBR-Induced Coherency Breakdown}

Classical WAPC coherency grouping (IEEE~1110, WECC stability criteria) clusters
generators by \emph{power angle synchronism}: machines whose rotor angles swing
together in response to a disturbance form one coherent group.  This grouping
determines which relays are coordinated together.

IBRs have no rotor angle.  Their ``coherency'' is defined by the PLL tracking
bandwidth and the reactive current injection priority: two IBRs at opposite ends
of a feeder may respond to the same voltage dip in diametrically opposite ways
(one injects reactive, the other active), creating an \emph{anti-coherent} pair.

Published WAPC designs (Apostolopoulos~2020, Sun~2022) apply SG coherency
algorithms to IBR-dominated networks, producing incorrect relay groups.  The SAMBP
WAPC (Ch~7) uses the EKF-estimated $\kibr$ as a coherency proxy, grouping relays
by their IBR penetration level rather than by power angle synchronism.

\subsection{G6 --- Adversarial ML Attacks on Protection Relays Unmitigated}

The adversarial ML literature (Goodfellow~2015, Madry~2018) demonstrates that
small crafted perturbations can cause neural network classifiers to misclassify
with high confidence.  In protection, this means:
\begin{itemize}[nosep]
  \item A $\leq 5\%$ perturbation in a CT secondary current can cause an
        unconstrained NN to classify a genuine fault as normal operation (miss trip).
  \item A $\leq 3\%$ perturbation can cause classification of load current as
        a fault (false trip).
\end{itemize}

The protection literature has not published a mitigation for this attack class.
The SAMBP response (Ch~6 and Ch~8) is two-pronged:
(a) the PINN loss function embeds KVL/KCL constraints, so any perturbation that
causes misclassification must also violate the power flow equations --- raising
the minimum attack energy; and (b) HMAC-SHA256 on GOOSE prevents the attacker
from injecting false measurements via the communication network.

%% ── 3. The Six Objectives: Verification Criteria ────────────────────────────
\section{The Six Objectives: Measurable Success Criteria}

\begin{keybox}{How Each Objective Is Verified}
The TSA will ask: \emph{how do you know you have achieved each objective?}

\begin{center}
\small
\begin{tabular}{@{}lp{4.5cm}p{5cm}@{}}
\toprule
\textbf{Obj.} & \textbf{Statement} & \textbf{Success Criterion (quantified)} \\
\midrule
O1 & Unified SAMBP LM framework
   & Convergence $R^2 > 0.999$ in $<$20 iterations for all models;
     Theorem~2.1 convergence proof \\[4pt]
O2 & Self-adaptive protection modules for 87L, SyncOC, 87T, 21, microgrid
   & TPR = 1.000; FPR = 0.000; avg.\ trip time $\leq 80$~ms across all elements \\[4pt]
O3 & Non-blocking EKF digital twin
   & $\hat{\kibr}$ RMSE $\leq 0.024$~pu; EKF background thread latency $\leq 2$~ms;
     relay decision not blocked \\[4pt]
O4 & Physics-informed ML for fault classification
   & Classification accuracy $>$95\% under $\varepsilon = 5\%$ CT perturbation;
     KVL/KCL residual $< 10^{-4}$~pu \\[4pt]
O5 & Wide-area coordinator with IBR coherency
   & Fault localisation accuracy 100\%; N-1 verified; settings dispatch $<$100~ms \\[4pt]
O6 & Cybersecurity layer
   & Spoofing success $< 0.001\%$; GOOSE latency overhead $< 1$~ms;
     risk reduction $>$50\% per IEC~62351 matrix \\
\bottomrule
\end{tabular}
\end{center}
\end{keybox}

%% ── 4. Prior Art: Three Competing Approaches ─────────────────────────────────
\section{Prior Art Analysis: Three Competing Approaches and Their Limits}

\begin{archbox}{How SAMBP Differs from Existing Approaches}

\paragraph{Approach A --- Legacy Relay Retuning (Hooshyar~2019, Sadeghi~2021):}
\textit{Premise:} Compute new fixed settings that are ``better'' for an IBR-heavy
network; apply them during commissioning or scheduled maintenance.\\
\textit{Limitation:} The setting is still fixed.  As IBR fleet composition
changes between maintenance cycles (seasonal wind variation, solar outages,
new capacity additions), the setting becomes stale.  At $\kibr=0.1$ vs.\
$\kibr=0.9$, the optimal bias slope differs by $4\times$; no single fixed
setting covers both.\\
\textit{SAMBP difference:} Real-time LM estimation tracks $\kibr$ continuously;
settings update every fault cycle.

\paragraph{Approach B --- Communication-Assisted Schemes (Gers~2011, Zhang~2020):}
\textit{Premise:} Use optical-fibre 87L (full differential) or POTT to
eliminate the need for fault current magnitude thresholds.\\
\textit{Limitation:} These schemes \emph{still have} a bias slope $\km$.
At $\kibr < 0.2$, the fixed $\km = 0.30$ causes a missed trip even with
full communication (as demonstrated in Ch~3 TR-08 results).  Communication
latency under cyberattack or fibre cut disables the scheme entirely.\\
\textit{SAMBP difference:} Adaptive $\km^*(\kibr)$ works at all $\kibr$
values; Stage-2 gate handles latency gracefully.

\paragraph{Approach C --- Pure Machine Learning (Jamali~2020, Liao~2022):}
\textit{Premise:} Train a deep neural network on simulation data covering
all fault types and IBR penetration levels; deploy as the relay decision function.\\
\textit{Limitation:} (a) Produces physics-impossible outputs (e.g., classifying
a fault where KVL is violated by the prediction); (b) generalisation to
out-of-distribution IBR types (e.g., DFIG crowbar activation) is poor;
(c) adversarial perturbation vulnerability is unaddressed in every published
protection ML paper.\\
\textit{SAMBP difference:} PINN loss embeds KVL/KCL; the LM model is physics-first;
adversarial hardening via HMAC + physics constraints.
\end{archbox}

%% ── 4b. Thesis Table 1.1 --- Verbatim ───────────────────────────────────────
\section{Thesis Table 1.1: Summary of Original Contributions (Authoritative)}

\begin{resultbox}{Table 1.1 from the Thesis (Chapter and Result Numbers Are Final)}
\begin{center}
\small
\begin{tabular}{@{}clp{8.5cm}c@{}}
\toprule
\textbf{ID} & \textbf{Ch} & \textbf{Contribution} & \textbf{Status} \\
\midrule
C1 & 3
  & Physics-constrained adaptive line differential (87L) protection with
    IBR-aware threshold $k_m^*(\kibr)$, achieving TPR\,=\,1.000 and zero
    false trips across the full IBR penetration range
  & [P] \\[4pt]
C2 & 4
  & Levenberg--Marquardt inverse-estimation framework for synchronous
    generator over-current protection, with bounded adaptive update law
    and four-component confidence gating ($R^2 > 0.999$)
  & [P] \\[4pt]
C3 & 6
  & Non-blocking EKF digital twin for real-time joint estimation of IBR
    penetration ratio $\kibr$, line impedance, and IBR admittance from
    IED phasor measurements (RMSE\,$\leq$\,0.024~pu)
  & [P] \\[4pt]
C4 & 7
  & Wide-area protection coordinator with PMU-based fault localisation,
    N-1 security verification, and automatic CTI-graded relay settings
    dispatch via IEC~61850 GOOSE
  & [P] \\
\bottomrule
\end{tabular}
\end{center}
{\small [P] = validated with full numerical simulation results.}

\medskip
\textbf{Correct chapter mapping for all cross-references:}
Ch~2 = Mathematical Foundations (LM framework) $\cdot$
Ch~3 = 87L Line Differential (C1) $\cdot$
Ch~4 = SyncOC Generator OC (C2) $\cdot$
Ch~5 = Adaptive Distance Protection $\cdot$
Ch~6 = EKF Digital Twin (C3) $\cdot$
Ch~7 = WAPC (C4) $\cdot$
Ch~8 = Cybersecurity
\end{resultbox}

%% ── 5. Novelty Claims: The Three Firsts ──────────────────────────────────────
\section{Novelty Claims: What SAMBP Does That No Prior Work Has Done}

\begin{resultbox}{Three Novelty Claims}
The thesis positions SAMBP as the \textbf{first work} to do all three of:

\begin{enumerate}[nosep]
  \item \textbf{Treat adaptive protection as a unified inverse estimation problem}
        with explicit model identification across all element types (87L, OC, 87T,
        distance).  Prior work treats each element in isolation with different
        estimation methods.  The SAMBP LM framework provides one solver,
        one confidence gate ($\gamma \geq \gammath$), and one convergence proof
        (Theorem~2.1) for all elements.

  \item \textbf{Integrate local relay adaptation, digital-twin state estimation,
        and centralised coordination} into a single closed-loop self-adaptive system.
        Prior works on adaptive protection, digital twins, and WAPC are published
        independently with no shared interface specification.  SAMBP defines the
        three-tier architecture (Tier~1 EKF oracle $\to$ Tier~2 protection
        modules $\to$ Tier~3 WAPC) with formal data contracts between tiers.

  \item \textbf{Provide formal confidence-gating and bounded update laws} that
        guarantee relay \emph{safety} under estimation uncertainty.  Existing
        adaptive relay papers show performance improvement without proving that a
        wrong estimate cannot push settings to an unsafe value.  The bounded update
        (Proposition~2.1) and Stage-2 gate provide this guarantee.
\end{enumerate}
\end{resultbox}

%% ── 6. Gap Mapping to Thesis Chapters ────────────────────────────────────────
\section{Gap and Objective Mapping to Thesis Structure}

\begin{center}
\small
\begin{tabularx}{\linewidth}{@{}llllX@{}}
\toprule
\textbf{Gap} & \textbf{Obj.} & \textbf{Ch} & \textbf{Contribution} & \textbf{Key Result} \\
\midrule
G1, G3 & O1 & 2   & Unified LM inverse-estimation framework & Theorem~2.1; $R^2 > 0.999$; $\leq 20$ iter \\
G1, G2 & O2 & 3--5 & Self-adaptive protection modules (87L C3, SyncOC C4, Dist C5) & TPR\,=\,1.000; FPR\,=\,0.000 \\
G4     & O3 & 6   & Non-blocking EKF digital twin (C3)      & RMSE\,$\leq$\,0.024~pu; 2~ms thread latency \\
G6     & O4 & 6   & Physics-informed NN (PINN) classifier   & $>95\%$ accuracy at $\varepsilon=5\%$ \\
G5     & O5 & 7   & WAPC with IBR-aware coherency (C4)      & 100\% localisation; 54.9\% risk reduction \\
G6     & O6 & 8   & Cybersecurity: HMAC + adversarial PINN  & Spoofing $<0.001\%$; overhead 0.72~ms \\
\bottomrule
\end{tabularx}
\end{center}

%% ── 7. Cross-Thesis Connections ──────────────────────────────────────────────
\section{Cross-Thesis Connections}

Every subsequent slide in the TSA presentation fills in one row of the table
above.  The logical flow is:
\begin{itemize}[nosep]
  \item \textbf{Slide~5} (Ch~2 LM framework) $\leftarrow$ closes G3 and G4.
  \item \textbf{Slide~6} (Ch~3 IBR characterisation) $\leftarrow$ quantifies G1
        and G2 with the failure map table.
  \item \textbf{Slides~7,~8} (Ch~3, C1: 87L) $\leftarrow$ delivers O2 for 87L.
  \item \textbf{Slide~9} (Ch~4, C2: SyncOC) $\leftarrow$ delivers O2 for OC.
  \item \textbf{Slide~11} (Ch~5: Distance) $\leftarrow$ delivers O2 for distance.
  \item \textbf{Slide~4} (Three-Tier Architecture) $\leftarrow$ shows O1--O6 integrated.
\end{itemize}

Slide~3 (this study note) is the \textbf{gateway} slide: the TSA should have it
clear in mind when evaluating every technical slide that follows.  If a committee
member challenges the novelty of any result, the answer traces back to one of the
six gaps stated here.

%% ── 8. Key Reference Points ──────────────────────────────────────────────────
\section{Key Reference Points for Literature Defence}

\begin{formulabox}{Standards and Papers to Name-Drop}
\begin{center}
\small
\begin{tabular}{@{}lp{9cm}@{}}
\toprule
\textbf{Reference} & \textbf{What It Says and How SAMBP Differs} \\
\midrule
IEC 60255-151       & Defines fixed-curve OC protection; no provision for real-time pickup adaptation \\
IEEE C37.112        & Standardises IDMT relay characteristics; no IBR model \\
Hooshyar (2019)     & IBR feeder OC retuning; offline process; no LM estimator \\
Sadeghi (2021)      & Distance relay reach correction for IBR; element-specific table; no unified framework \\
Adamiak (2018)      & Adaptive 87L for HVDC; DC-source specific; not GFL/GFM \\
Goodfellow (2015)   & Adversarial examples in ML; SAMBP applies this to protection relays and mitigates it \\
Jamali (2020)       & DNN fault classifier for IBR networks; no physics constraints; adversarial not addressed \\
IEEE PES DT TF (2022) & Digital twin survey; no bounded update law; no non-blocking guarantee \\
Taul (2020)         & GFM protection-control coupling analysis; no relay algorithm proposed \\
\bottomrule
\end{tabular}
\end{center}
\end{formulabox}

%% ── 9. Speaker Script ────────────────────────────────────────────────────────
\section{Timed Speaker Script (3 Minutes)}

\begin{speakbox}
\textbf{[0:00--0:30] Setting the Scene}\\
``Slide~3 identifies the six specific gaps that this thesis fills.  I grouped them
into three clusters.  The first two are standards and modelling gaps: no relay
standard in IEC~60255 or IEEE~C37 accounts for real-time IBR admittance, and the
protection-control coupling unique to grid-forming IBRs has not been addressed in
any relay design.  These two gaps mean that even if a protection engineer wanted
to build an adaptive relay today, there are no published methods to follow.''

\medskip
\textbf{[0:30--1:10] Algorithmic and DT Gaps}\\
``The next two are algorithmic gaps.  There is no unified inverse-estimation
framework: the OC literature has one method, the 87L literature has a different
method, and distance has a third.  None of them share a common solver, confidence
gate, or convergence proof.  The SAMBP Chapter~2 provides that unified framework.
The digital twin gap is more specific: every published protection digital twin
either runs offline or runs blocking in the relay's main thread.  Neither has a
formal proof that settings stay within safe bounds during a bad estimate.  My
bounded update law and Stage-2 gate fill this gap.''

\medskip
\textbf{[1:10--1:50] System and Security Gaps}\\
``The last two are system-level gaps.  Wide-area protection coordinators today
group relays by power-angle coherency --- a concept that only applies to synchronous
generators.  IBRs have no rotor angle.  My WAPC groups relays by IBR penetration
ratio instead, using the EKF estimate as the coherency proxy.  And the adversarial
ML gap: it has been shown in computer vision that a 5\% perturbation in the input
can fool a neural network.  For a protection relay, that means a 5\% injected CT
error causes a misclassification.  No protection ML paper has addressed this.
My PINN loss function and HMAC countermeasure address it.''

\medskip
\textbf{[1:50--2:30] The Six Objectives}\\
``The six research objectives on the right close the six gaps one-to-one.  The
unified SAMBP framework closes G1 and G3.  Self-adaptive modules close G1 and G2
across all five protection element types.  The non-blocking EKF digital twin closes
G4.  Physics-informed ML closes G6's classifier aspect.  The wide-area coordinator
closes G5.  And the cybersecurity layer closes G6's communication aspect.  Each
objective has a quantified success criterion: TPR~=~1.000, FPR~=~0.000, RMSE
$\leq$~0.024~pu, spoofing below 0.001\%.''

\medskip
\textbf{[2:30--3:00] Novelty Position}\\
``I want to be explicit about what is new.  This is the first work to treat
adaptive protection as a single unified inverse estimation problem across all
relay types.  The first to integrate local adaptation, digital twin, and
centralised coordination into one closed-loop system.  And the first to provide
formal bounded-update and confidence-gate guarantees that prove safety under
estimation uncertainty.  Prior adaptive relay work improves performance; this
thesis guarantees it.''
\end{speakbox}

%% ── 10. Anticipated Q&A ──────────────────────────────────────────────────────
\section{Anticipated Q\&A --- Five Likely Committee Questions}

\begin{keybox}{Q1: How do you know no relay standard accounts for real-time $\kibr$? What about IEC 61968 or IEC 62351?}
\textbf{Answer:}  IEC~61968 is a CIM (Common Information Model) standard for utility
data exchange --- it defines data schemas, not relay algorithms.  IEC~62351 is the
cybersecurity series for IEC~61850 communications --- it also does not define relay
algorithms.  The relay algorithm standards are IEC~60255-151 (overcurrent), IEC~60255-121
(distance), and IEEE~C37.112 (inverse time); none of these define a mechanism for
real-time relay threshold computation from an online $\kibr$ estimate.  I have reviewed
all three standards in Chapter~1's literature review.  The closest is IEC~60255-121
Annex~B, which allows for ``special characteristics'' by mutual agreement between
manufacturer and user --- this is essentially a blank cheque that acknowledges the
standard is insufficient, which is exactly the gap.
\end{keybox}

\begin{keybox}{Q2: What specifically is ``IBR-induced coherency breakdown'' in G5?}
\textbf{Answer:}  Classical coherency in power systems means that during a disturbance,
a group of synchronous generators swing together in power angle --- their rotor speeds
stay nearly identical, so they form one ``coherent group'' for protection coordination.
When you replace SGs with IBRs, coherency breaks down in two ways: (a) IBRs have no
rotor --- they have no inertia and no swing equation; (b) IBRs actively inject current
in a direction commanded by the LVRT control law, not by physics.  Two IBRs at
opposite ends of a feeder can inject reactive current in \emph{opposite} directions
simultaneously, creating an anti-coherent pair.  The WAPC must know which relays are
protecting IBR-dominated lines to apply different coordination criteria.  The SAMBP
WAPC uses $\hat{\kibr}$ from the EKF as a coherency proxy: relays with similar $\hat{\kibr}$
values are grouped together for CTI coordination.
\end{keybox}

\begin{keybox}{Q3: ``Protection-control coupling of GFM IBRs unstudied'' seems strong. Are there no papers on this?}
\textbf{Answer:}  The analysis of GFM protection challenges exists in the power
electronics literature (Taul~2020, Lasseter~2020).  What is \emph{unstudied} is a
relay algorithm that explicitly accounts for the coupling.  Taul~2020 shows that
changing the virtual impedance $Z_v$ to increase fault current for relay visibility
degrades the GFM's voltage regulation stability --- but proposes no relay solution.
The SAMBP microgrid module (Ch~8) is the first relay algorithm that leaves $Z_v$
untouched and adapts the relay settings instead, operating entirely within the GFM's
existing control constraints.  So the gap is precisely stated: the algorithm has not
been published, even though the problem has been identified.
\end{keybox}

\begin{keybox}{Q4: You have 6 gaps and 6 objectives but the thesis has 4 principal contributions (C1--C4). How do these reconcile?}
\textbf{Answer:}  The six objectives on the slide are \emph{granular}; the four
principal contributions (C1--C4) are \emph{consolidated} for thesis chapter organisation:
From Table~1.1 of the thesis: C1 (Ch~3) is the adaptive 87L line differential achieving
TPR\,=\,1.000; C2 (Ch~4) is the LM inverse-estimation framework for SyncOC with
$R^2>0.999$; C3 (Ch~6) is the non-blocking EKF digital twin with RMSE\,$\leq$\,0.024~pu;
C4 (Ch~7) is the WAPC with PMU fault localisation and CTI-graded settings dispatch.
All four are validated~[P].  The six-objective structure on the slide unpacks the
four chapters into finer deliverables (e.g., O4~=~PINN and O6~=~cybersecurity are both
inside Ch~6/Ch~8 respectively), making the novelty mapping explicit for the committee.
\end{keybox}

\begin{keybox}{Q5: Are there commercial products (ABB, Siemens, GE) that already do adaptive protection? How is your work different?}
\textbf{Answer:}  Yes --- ABB REL~670, Siemens SIPROTEC~7SA87, and GE Multilin~D90
all claim ``adaptive'' protection.  Their adaptation mechanism is \emph{setting groups}:
the relay engineer pre-configures 2--8 groups of fixed settings, each for a different
expected network condition, and an external trigger (SCADA topology message or digital
input) switches between groups.  This is \emph{not} real-time inverse estimation: the
groups are computed offline, and the relay does not identify any parameter from its
own CT measurements.  It is essentially ``manual look-up table'' adaptation.  SAMBP
is fundamentally different: the relay computes its own setting from a real-time
physical model fit.  The analogy: commercial ``adaptive'' relays are like a photographer
choosing between pre-set exposure modes (portrait, sport, night); SAMBP is like a
camera that measures the actual scene lighting and computes the correct exposure.
\end{keybox}

\vspace{1em}
\begin{tcolorbox}[colback=iitmnavy!5,colframe=iitmnavy,arc=3pt,boxrule=1pt]
\small\centering
\textbf{Summary for the TSA Committee}\\[3pt]
Six identified gaps (two standards, two algorithmic, two system-level) map to six
research objectives, each with a quantified success criterion.  SAMBP's novelty
rests on three claims: unified inverse-estimation across all relay types, closed-loop
integration of adaptation and coordination, and formal safety guarantees via bounded
updates and confidence gating.  Every technical slide that follows delivers one of
these six objectives.
\end{tcolorbox}

\end{document}
